\section{Joint interpretation}

The empirical contribution is a common prescription that connects a
galactic acceleration scale, a retained mass distribution, tidal evolution,
and cosmological energy transfer. The calculations use the same
microscopic entropy and conversion fraction in every sector, so a change
of either coefficient moves several observables at once.

Against MOND, the retained stock supplies the cluster mass, the LMC mass
inferred from its stream, the vertical gravity of the NGC~628 old disk,
and a fit to the acoustic spectra. MOND fits the isolated-lens profiles
better, predicted the Crater II dispersion in advance, shares the
rotation-curve law, and fits SPARC with a lower $\chi^2$ when its $a_0$ is
fitted. Against \LCDM, the stock model matches the cosmological
fit with two fewer parameters, predicts $\Omega_\Lambda$, $H_0$ and the
matter density, fixes $a_0$, and reaches the LMC and NGC~628 values with
no halo freedom. \LCDM\ fits Antlia 2 and, at low concentration, the
isolated-lens profiles more closely. Table~\ref{tab:headtohead} gives the
values and Section~\ref{sec:open} the open items.

\subsection{The derivation has identifiable physical commitments}

\begin{table}[htbp]
\small
\begin{tabularx}{\textwidth}{P{.24\textwidth}YY}
\toprule
Step & Result supplied & Physical commitment or outstanding test \\
\midrule
Finite face and cell ensemble
& $16=7+9$; 1680 cell states; $g=7.4198000$ and $\eta_*=0.02986684$.
& The spin, matching vertex, occupancy rule, closure invariant, and
evidence measure define the ensemble. \\
Horizon and transverse response
& $\epsilon=g/(2\pi)^2$; $a_0=\epsilon cH$; Bose response.
& The tilt-mode pendulum length $\ell=g\,c/H$ in the horizon bath connects the
count to acceleration and fixes the mode energy. \\
Retained galactic stock
& An inertial positive mass component, inherited by carriers and
transported by collisionless evolution.
& Assembly budgets, carrier assignment, profile preparation, and
relaxation must be specified independently of the target observables. \\
Cosmological conversion
& $\eta=\eta_*$ and an initial abundance near $\epsilon^{-1}$.
& Common energy conversion, the selected pre-acoustic formation map
(whose microscopic vertex remains open), and the absolute vacuum
matching condition enter the physical identification. \\
Assembly and mixing
& A cumulative release function and a weighted vacuum integral.
& Retention gates, progenitor labels, and transport set the normalization;
the current expansion fit determines one mixing parameter. \\
Geometry and perturbations
& Reduced Newtonian equations and no-slip identity proved from the
specified action; carrier equations proved on their regular branch.
& The continuum embedding, constrained director dynamics, and the
physical matching of the effective actions remain separate requirements. \\
\bottomrule
\end{tabularx}
\caption{The logical connection between the microscopic model and the
observables. An exact finite calculation can coexist with conditional
physical identifications and an incomplete continuum derivation.}
\end{table}

This chain explains why the empirical results bear on the microstructure:
the model does not choose a separate acceleration coefficient for the LMC
and the redshift sample, or a separate release fraction for each
cosmological data combination. It also identifies which steps the data
cannot determine by themselves. A fit to an effective gravitational field
does not uniquely identify the microscopic state space that produced it.

\subsection{Which conventional assumptions are challenged}

For the LMC, the comparison challenges the applicability of an
instantaneous external-field suppression to a previously established
gravitating stock. A retained distribution continues to contribute after
the ambient field changes. The tested QUMOND prescriptions and the stock
model make different mass predictions with the same stated baryons.
Joint stream and reflex modeling can test that distinction more directly
than comparing scalar enclosed-mass summaries.

For Crater II, the issue is the central structure and stripping response
of the excess gravitating mass. The tested capacity profile loses enough
support to produce a cold, extended remnant under the adopted orbit and
initial conditions. Standard-concentration NFW halos under the tested
abundance prior perform poorly in that comparison, and MOND reaches the
same dispersion through the external field. Low-concentration
and cored dark-matter alternatives require their own matched tests.
Assembly history already matters in conventional dark-matter dynamics;
history dependence alone cannot distinguish the two frameworks.

For the acceleration scale, the proposed change is a cosmological
normalization $a_0\propto H$. Its observed value at one epoch would cease
to be an independent constant. The published intermediate-redshift
relation is consistent with the fixed reference curve under the stated
projection. A direct galaxy-level likelihood would test both its curvature
and its scatter.

Irreversible source-identity loss generates vacuum energy. The microscopic
conversion and formation abundance tie the release history to the CMB
and distance fits. The halo occupation of feedback-resolved simulations sets the retention
gate with no input adjusted to the vacuum density and gives
$\Omega_\Lambda=0.709\pm0.02$ against the required $0.707$. Gates computed
without feedback give $0.73$ from the hydrogen cooling criterion and up to
$0.77$ from condensation in the archived run
(Section~\ref{sec:coolinggate}). The frozen history, with the reionization redshift taken from the fitted optical depth, fits the cosmological data with four parameters at $\Delta\chi^2=+1.07$ against six-parameter \LCDM. A measured cold-gas retention fraction sharpens this normalization test. The current evidence supports these mechanisms as
testable alternatives; it does not establish the elimination of all
dark-matter models, all MOND formulations, or the absolute cosmological
constant problem.

\appendix
\section{Provenance and numerical verification}

\subsection{Scope of the verification}

The integrated Lean 4.35.0 library contains 160 modules and 2,937 audited theorem
constants, including its legacy sublibrary. The manuscript ledger records
874 statements as proved, 90 in part, 153 as physical premises, 64 as
data comparisons, and 19 as simulation outputs. Its other entries
distinguish numerical calculations, definitions, and statements awaiting
formalization. The saved audit reports only the standard Lean axioms
and no admitted proofs. The counts include mathematical helper results.
\source{T4}

For this paper, the relevant proofs cover the finite ensemble and closure
root, the restricted freedom in the marked-register identification,
positive finite allocation, and radial relaxation with preserved angular
marginals. First variation with compactly supported tests gives the
reduced two-potential and Newton equations. Decay at infinity gives
$\Phi=\Psi$. The carrier-action proof assumes a smooth interpolating
function; the adopted singular derivative at zero baryonic field limits
its direct classical application. The constrained director equation,
full cosmological perturbations, and observational likelihoods retain
their separate analytic or numerical status.

The numerical review independently evaluates the finite ensemble,
redshift projection, LMC mass profile and QUMOND flux, archived Crater II
likelihoods, and a reference release history. The archived verification
record in \path{evidence/legacy_review} uses the unrounded acceleration
coefficient and the higher-accuracy matched cosmology endpoints.
The associated figures use those records. The current
\path{current_verification.json} records the saved three-theory comparisons
and manuscript integrity checks. The archived release-history and
hydrodynamic panels retain their original linear-depletion calibration.

The dated research records include the exponential-allocation cosmological
optimization, finite-work and fluid tests, merger controls, observed cluster
reductions, and outlier audits. They were not rerun to typeset this paper.
Six extension modules contribute 119 explicit theorem declarations.
The integrated 4.35 build, source scan, axiom audit, and ledger citation
check pass; the six modules also pass standalone Lean/mathlib 4.19 checks.
The existing claim ledger retains its snapshot scope. \source{T5} Source inspection and saved
audits establish which calculation supports each statement. Numerical
convergence, selection effects, and physical matching are assessed with
the relevant empirical result. SHA-256 records preserve the identities
of the source and result snapshots.

\clearpage
\subsection{Local source key}

Paths below are relative to the supplied \path{Sol and Fable} workspace.
The source package includes the document, figures, numerical review, and
verification script; the original simulation archives remain in that
workspace.

\small
\begin{longtable}{P{.065\textwidth}P{.88\textwidth}}
\toprule
Key & Source and use \\
\midrule
T1 & \path{01 Current Papers/Full/Entropic Scalar EFT.tex}. Current microscopic
construction, effective galaxy dynamics, release accounting, and
covariant closure; especially sections 5--6, 15, 20, N.8, and O.3--O.15.
The numerical review records its source snapshot by hash. \\
T2 & \path{02 Research/Spacetime consolidation 2026-10-16/MICROSTRUCTURE_TO_SPACETIME.md}.
Finite geometry, interaction, matter, and clock
construction. \\
T3 & \path{01 Current Papers/Microstructure CDT/MICROSTRUCTURE_CDT.tex}.
Selected CDT realization, normalization conventions, and finite
geometry--matter--clock tests, with conventions aligned to the current
Entropic Scalar EFT manuscript. \\
T4 & \href{https://github.com/jakechinitz/Entropic-EFT-Lean/tree/cfcfa373eaf55adac6a2a0ddeb859766c985f71c}{Entropic-EFT-Lean, 29 October 2026 snapshot}.
Local library: \path{02 Research/Unified closure 2026-10-21/formalization}.
Claim ledger, theorem sources, and saved axiom/citation audits. \\
G1 & \path{04 Notes and References/Technical Notes/history_gravity_reproduction_2026-09-27.md}.
Analytic galaxy checks and reproducible provenance for the stock model. \\
G2 & \path{02 Research/LMC and satellites checks/qumond_flux.py} and
\path{02 Research/LMC and satellites checks/output.txt}.
External-field QUMOND surface-flux comparison. \\
G3 & \path{02 Research/Crater II and Antlia 2 vs LCDM/PREREGISTRATION.md},
\path{PREREGISTRATION.sha256}, and \path{analyse_grid.py} in that folder.
Core observation likelihood, priors, and hash identity. \\
G4 & \path{02 Research/Crater II and Antlia 2 vs LCDM/RESULTS_with_LMC.md},
\path{analyse_lmc.py}, and \path{lmc_results.zip} in that folder.
Current moving-LMC trajectories and evidence. \\
G5 & \path{04 Notes and References/Technical Notes/carrier_hierarchy_calcs_2026-09-26.md}.
The active-carrier inheritance calculation for tidal dwarfs. \\
G6 & \path{02 Research/Crater II and Antlia 2 vs LCDM/bf_robustness.txt} and
\path{antlia2_ji_mock/mock_approx_analysis.txt} in that folder.
Model-error sensitivity and approximate observational mocks. \\
G7 & \path{02 Research/Stock adjustment resolution 2026-10-21/START HERE.txt};
\path{02 Research/Halo replication audit 2026-10-21}.
Radial adjustment, retained angular marginals, and fixed-target halo tests. \\
G8 & \path{02 Research/Galaxy sector resolution 2026-10-24} and its
\path{Repairs - 3D contacts and exact minimiser} folder. Selection
constraints, exact constrained minimiser, three-dimensional contact runs,
lensing bound, and independent audits. \\
G9 & \path{02 Research/Contact rule physical basis 2026-10-24}. Bose
relative-entropy bias, equipartition weight, zero-temperature detailed
balance, Bose-functional runs, and the Lean module
\path{Entropic/ContactSelection.lean}. \\
G10 & \path{02 Research/Bose contact observed clusters 2026-10-24}. Forty
contact runs on seven X-COP clusters without a cluster gate, at
galaxy-scaled and fixed work allowances, with numerical audit. \\
G11 & \path{02 Research/Cluster switch derivation 2026-10-24}. Record-dominated
supports and cell energies of the X-COP clusters, the cell reservoir
$E_{\rm cell}/M_{\rm stock}\to V_f^2$, contact runs under four rules, and gated
galaxy runs. \\
G12 & \path{02 Research/Highest-leverage tensions 2026-10-28}: \path{lensing/},
\path{dwarfs/}, \path{ELIMINATION.md} and \path{REVIEW_CORRECTIONS.md}.
Isolated-lens population forward model with support-limited stock, and
the satellite support assignment for the classical dwarfs, with
independent review. \\
C1 & \path{02 Research/Full Planck tied abundance/model.py},
\path{RESULTS.md}, and \path{fits_lenscorr.jsonl} in that folder;
\path{02 Research/Dark energy prediction tests/common.py}.
Matched likelihood, data handling, Boltzmann mapping, and reference
\LCDM\ result. \\
C2 & \path{02 Research/Release zoom/RESULTS.md}, \path{zoom_fits.jsonl},
\path{zoom_fit.py}, \path{o12_release.py}, \path{Frel_table.npy},
\path{FZ_ref.npy}, and \path{kmix_summary.txt} in that folder.
Archived PPF release fit, reference history, and hydrodynamic mixing test. \\
C7 & \path{02 Research/Dark energy normalization 2026-10-30}: \path{cooling_gate.py},
\path{gate_from_run.py}, \path{irel_gate.py}, their JSON outputs, and \path{NOTE.md}.
Retention gate from the cooling criterion and from the archived run, and
the release integral, $h$ and $\Omega_\Lambda$ against the gate. \\
C8 & \path{02 Research/Retention gate from feedback simulations 2026-10-31}: \path{gate_mass.py},
\path{grid_mass.py}, \path{mixtures.py}, \path{renvar.py}, their JSON outputs, and \path{NOTE.md}.
Release integral with the feedback-resolved occupation as the retention gate, and its uncertainty;
\path{frozen_likelihood/} holds the frozen-history fits. \\
C3 & \path{02 Research/Formation operator check/RESULTS.md}.
Formation abundance and the $5.333$ mass convention. \\
C4 & \path{02 Research/High leverage closure 2026-10-21/fit_comparison.json},
\path{fit_physical_geodesic/result.json}, \path{fit_physical_lcdm/result.json},
\path{cosmology_checks.json}, and \path{growth_independent_check.json} in that folder.
Direct clock-fluid implementation, matched linear-lensing fits, and numerical checks. \\

T5 & \path{02 Research/Manuscript cluster and dynamics integration 2026-10-23}.
Preserved source hashes, integration diffs, six-module manifest, compilation
checks and PDF validation. The six modules are included in the
\href{https://github.com/jakechinitz/Entropic-EFT-Lean/tree/a6552daf39daf83d5159fcbcda698d1801731d24}{checked extension snapshot}
under \path{Entropic/Extensions20261023}. \\
T6 & \path{02 Research/Empirical fifteen-test structure 2026-10-23}.
Fifteen-family coverage map, preserved source versions, editorial diffs,
and validation of the two-part presentation. \\
U1 & \path{02 Research/Unified empirical program 2026-10-23/execution_register.json}.
Fifteen-sector evidence register; S01--S15 identify the selected datasets,
saved results, source URLs, and remaining tests. \\
C5 & \path{02 Research/Unified empirical program 2026-10-23/results/cosmology_final.json}.
Exponential-allocation fit, matched reference, and saved observables. \\
C6 & \path{02 Research/Highest-leverage tensions 2026-10-28/cosmology}:
\path{harness_v2.py}, \path{run/}, and \path{REVIEW_CORRECTIONS.md} in the
parent folder. One-harness \LCDM\ and direct geodesic release fits across
supernova, neutrino and lensing settings, with independent review. \\
D1 & \path{02 Research/Connected theory closure 2026-10-23}.
Depletion action, unchanged linear domain, lag and noise bounds, and
conditional tensor calculation; 42 Lean declarations. \\
D2 & \path{02 Research/Effective dynamics and coupling 2026-10-23}.
Finite-work contact, warm fluid stress, compression and feedback checks;
36 Lean declarations. \\
CL1 & \path{02 Research/Conservative cluster merger 2026-10-23}.
Seventeen encounter/control runs, conservation and refinement records;
8 Lean declarations. \\
CL2 & \path{02 Research/Observed cluster tests 2026-10-23}.
X-COP and MENeaCS profiles, Bullet v1 targets, source hashes, and
15 numerical/data-handling checks. \\
CL3 & \path{02 Research/Cluster outlier closure 2026-10-23}.
Conditional pressure and transport audit, A644 positivity test, twelve
finite transport maps, 14 saved checks, and 22 Lean declarations. \\
X1 & \path{02 Research/Three-theory comparison 2026-11-08/sparc}. SPARC
under the stock law, MOND with three interpolating functions, and NFW halos
on the archived pipeline; prior and sample sensitivity of $a_0$. \\
X2 & \path{02 Research/Three-theory comparison 2026-11-08}: \path{clusters/},
\path{lmc_ngc628/} and \path{RESULTS.md}. MOND on the X-COP profiles,
abundance-matched \LCDM\ LMC masses and NGC~628 halo, and information
criteria. \\
X3 & \path{02 Research/Three-theory comparison 2026-11-08/kids}. MOND,
support-limited stock and NFW on the KiDS isolated-lens forward model with
matched nuisances. All three comparisons were rerun by an independent
audit. \\
\bottomrule
\end{longtable}
\normalsize

\subsection{Reproduction commands}

The source package includes the saved X1--X3 comparison outputs. From its
directory, the numerical and manuscript checks run with
\begin{verbatim}
python audit_saved_results.py
python figures/make_comparison_figures.py --output-dir figure_check
\end{verbatim}
The first command checks the saved scores, information-criterion arithmetic,
test-family coverage and references. The second regenerates the eight
comparison figures from the supplied outputs. Neither command fits a
likelihood or runs a simulation.


With Python, NumPy, SciPy, and Matplotlib installed, run
\begin{quote}\ttfamily
python reproduce\_review.py -{}-workspace "PATH/TO/Sol and Fable"
\end{quote}
The script evaluates archived inputs without launching simulation jobs.
The supplementary \path{refresh_verified_results.py} updates the
unrounded redshift arithmetic and imports the saved exponential-allocation
fit. It preserves prior outputs before writing. Its
\texttt{-{}-legacy-linear} option reproduces the earlier matched benchmark.
The broader calculations and source records are identified by U1 and C5.
Build the publication from
the workspace root with
\begin{quote}\ttfamily
python .organization/build\_pdfs.py -{}-paper empirical
\end{quote}
The build preserves previous PDFs and stores compiler files in the
workspace build directory. The source archive supplies the figures for
standalone typesetting.

The public dark-energy release and DES data repository are linked in the
references. The current physical mixing extension is identified by its
local archive and hashes; availability of the public reproduction bundle
does not by itself make every subsequent calculation public.

\section{Evidence register}
\label{app:register}

The register lists the fifteen test families recorded in the execution
register, with the present reach of each. The grouping records how each
family was registered; the text places each test by regime.

\begingroup\small
\setlength{\tabcolsep}{4pt}
\renewcommand{\arraystretch}{1.10}
\begin{longtable}{@{}p{.06\textwidth}p{.24\textwidth}p{.56\textwidth}p{.06\textwidth}@{}}
\caption{All fifteen test families in the evidence register. Placement
identifies the role of a test; the status column states its present reach.
The register is exploratory, with the existing outcomes known.}\label{tab:fifteen}\\
\toprule ID & Test & Present result and limit & See\\\midrule
\endfirsthead
\toprule ID & Test & Present result and limit & See\\\midrule
\endhead
\multicolumn{4}{@{}l}{\textbf{Families registered as distinctive tests}}\\[3pt]
S01 & High-redshift kinematics & Reference line projection agrees; galaxy catalogue residuals retain a high-redshift tension. & \ref{sec:redshift}\\
S02 & CMB--galaxy normalization & Exponential-fit $a_0=1.25883\times10^{-10}\,\mathrm{m\,s^{-2}}$; standard linear-depletion value $1.25005\times10^{-10}\,\mathrm{m\,s^{-2}}$; conditional mass and survival factors. & \ref{sec:crossnorm}\\
S10 & Binaries and Solar System & Cassini null benchmark is compatible; a Gaia catalogue likelihood remains open. & \ref{sec:binaries}\\
S11 & $G$ and lepton masses & Close selected-formula values; finite routing tests leave discovery bias unquantified. & \ref{sec:constants}\\
S13 & Milky Way, LMC, streams, tilt & LMC mass agreement and repeated tilt realizations; Galactic mass and a joint stream fit remain unresolved. & \ref{sec:lmc}, \ref{sec:halo}\\
S14 & Satellites and tidal dwarfs & Crater II favors stock within the tested grids; Antlia 2 and odds depend strongly on model error. & \ref{sec:crater}\\
S15 & Old-disk vertical gravity & NGC~628 gives $\chi^2=2.04$ for five dispersions; a broader tracer and force comparison is needed. & \ref{sec:olddisk}\\
\midrule
\multicolumn{4}{@{}l}{\textbf{Families registered as recovery of established phenomenology}}\\[3pt]
S03 & Local RAR and BTFR & Broad SPARC diagnostic and small BTFR scatter; the score is sensitive to velocity errors. & \ref{sec:localrar}\\
S04 & Galaxy weak lensing & Fifteen-bin diagonal score 10.06; reduction sensitivity and joint covariance remain. & \ref{sec:galaxylensing}\\
S05 & CMB acoustic peaks & Planck-based matched fit and ACT comparison recover similar quality to the reference cosmology. & \ref{sec:cmb}\\
S06 & Growth and Weyl power & ACT lensing score 16.21 over ten bins; DESI full-shape and nonlinear growth tests remain open. & \ref{sec:cmb}\\
S07 & Dark-energy history and distances & Matched BAO/SN fit with exponential allocation; the feedback-resolved retention occupation gives $\Omega_\Lambda=0.709\pm0.02$ against $0.707$ with no input adjusted to it, and the frozen history fits with four parameters at $\Delta\chi^2=+1.07$. & \ref{sec:distances}\\
S08 & Relaxed and colliding clusters & Positive stock targets and conserved encounters; outlier mechanisms and lens-map fits remain conditional. & \ref{sec:clusters}\\
S09 & Slip and $E_G$ & Transfer predictions are compatible with two published estimates; a joint window-convolved fit is needed. & \ref{sec:eg}\\
S12 & Waves and Lorentz tests & Luminal tensor branch is compatible with the polarization benchmark; operator-level Lorentz tests remain open. & \ref{sec:tensor}\\
\bottomrule
\end{longtable}
\endgroup
\source{U1; CL1--CL3; D1--D2; T6}

\section{Notation}
\small
\begin{tabularx}{\textwidth}{P{.23\textwidth}Y}
\toprule Symbol & Meaning \\
\midrule
$g$ & Cell sharing entropy in natural-log units. \\
$\epsilon=g/(4\pi^2)$ & Coefficient of the tilt-mode pendulum length over $(2\pi)^2$. \\
$\eta_*$ & Closure evidence coefficient; also the release fraction under
the common energy-conversion prescription. \\
$a_0$ & Characteristic galactic acceleration scale. \\
$g_b$ or $g_N$ & Newtonian acceleration produced by the assigned baryonic
source. \\
$n_B$ & Occupation of the adopted transverse mode. \\
$M_b$, $M_b^L$ & Present baryonic mass and ancestral Lagrangian baryonic
mass, respectively. \\
$C_R$ & Adopted retained-stock cap relative to present baryonic mass. \\
$R_h$ & Projected stellar half-light radius. \\
$\sigma_{\rm los}$ & Line-of-sight velocity dispersion, with sample and
gradient treatment specified for each comparison. \\
$Z$, $B$ & Prior-integrated likelihood and its ratio between two models. \\
$F_Z$ & Fraction that has first entered a baryon-retaining well. \\
$F_{\rm rel}$ & Cumulative fraction processed by the source-identity
release kernel. \\
$f$ & Mixing time in units of the formation halo dynamical time in the
pure-exponential fit. \\
$\rho_{c0}^{\rm wb}$ & Would-be present commitment density before release
losses, extrapolated as dust. \\
$I_{\rm rel}$ & Linear-benchmark release integral weighted by $a^{-3}$. \\
$I_{\eta,\rm rel}$ & Positive-allocation integral with the additional
survival factor $e^{-\eta_*F_{\rm rel}}$. \\
$\omega_i=\Omega_i h^2$ & Physical density parameter in the cosmological
fit. \\
\bottomrule
\end{tabularx}
\normalsize
